\documentclass[11pt,a4paper]{article}

\usepackage[margin=22mm,headheight=15pt]{geometry}
\usepackage[no-math]{fontspec}
\setmainfont{Noto Serif}
\setsansfont{Noto Sans}
\setmonofont[Scale=0.82]{DejaVu Sans Mono}
\usepackage{polyglossia}
\setdefaultlanguage{english}
\usepackage{amsmath,amssymb}
\usepackage{array,booktabs,graphicx,tabularx,xltabular,xcolor}
\usepackage{xurl}
\usepackage[unicode,colorlinks=true,linkcolor=blue!55!black,urlcolor=blue!55!black]{hyperref}
\usepackage{bookmark}
\usepackage{enumitem}
\setlist{itemsep=3pt,topsep=4pt,leftmargin=*}
\usepackage{titlesec}
\titleformat{\section}{\large\bfseries}{\thesection}{0.6em}{}
\titleformat{\subsection}{\normalsize\bfseries}{\thesubsection}{0.6em}{}
\setlength{\emergencystretch}{3em}

\title{MLC proof progress and the virtual-Molecule roadmap}
\author{Lean formalization status report}
\date{\today}

\begin{document}
\maketitle

\section{Verified state}

The repository contains a checked conditional reduction
\[
  \begin{aligned}
  &\texttt{Problem43PseudoSiegelAPrioriBoundsData}\\
  &\land\texttt{Problem44VirtualMoleculeData}\\
  &\Longrightarrow
  \texttt{LocallyConnectedSpace}(M).
  \end{aligned}
\]
The checked proof uses only \texttt{propext}, \texttt{Quot.sound}, and
\texttt{Classical.choice}. It contains no project \texttt{axiom},
\texttt{sorry}, or \texttt{admit}.

The currently proved foundation includes:

\begin{itemize}
  \item compactness and connectedness of the filled Julia set \(K_c\) for
        \(c\in M\);
  \item \(J_c\subseteq K_c\);
  \item a generic theorem producing an inclusion-minimal compact connected
        hull in a compact continuum;
  \item phase-indexed finite critical-orbit windows and partially invariant
        virtual hull data;
  \item transfer of a virtual separation estimate to the parameter-side
        separation estimate on connected nests;
  \item compactness of \(M\), the inclusion \(M\subseteq
        \overline{B(0,2)}\), and the certified interior inclusion
        \(B(0,1/4)\subseteq\operatorname{int}M\);
  \item a kernel-checked dyadic grid of depth 13 (with coarser squares of
        levels 10 and 11 in the shells of hyperbolic components) with
        first-order models of the critical orbit: squares proved outside \(M\)
        or inside \(M\) are removed, giving
        \(\mu^*(M\setminus\operatorname{int}M)\leq 2\cdot556432512/4^{16}
        \approx 0.259\), or \(2.1\%\) of the disk envelope.
\end{itemize}

The two remaining root obligations are genuine mathematical existence
statements:

\begin{enumerate}
  \item pseudo-Siegel a priori realization for all remaining unbounded
        satellite quadratic-like cases;
  \item virtual-Molecule completion at every uncovered boundary parameter by
        finite-return, connected virtual Julia hulls with a divergent a priori
        separation budget.
\end{enumerate}

Both are decomposed into statements in \texttt{Mlc/ProgrammeRoot.lean}: the
structure \texttt{ProgrammeInput} has one field per open statement of the programme,
sixteen in all (a transfer statement, six published cases, three statements of
Problem~4.3 and six of Problem~4.4), and \texttt{mlc\_of\_programme} proves in Lean
that they give the root input. The skeleton of the argument of \S4.1--\S4.3 is checked: a cap on the
jumps of the degeneration of the renormalization levels (a jump by a fixed factor
lands below a fixed ceiling) and a dichotomy with an error at most linear in the
period give a priori bounds for every tower
(\texttt{aPrioriBounds\_of\_jumpCap}). Given the dichotomy, the cap is also
necessary for the a priori bounds (\texttt{aPrioriBounds\_iff\_jumpCap}), so it is
the weakest hypothesis of its shape, and the Teichm\"uller contraction caps the
jumps at the levels of small relative period, so that the open content is the
tail at large relative period. Both ingredients are necessary
(\texttt{uniformities\_necessary}): for every unbounded \(h\), a tower with
unbounded relative periods satisfies the Teichm\"uller contraction, a uniform
growth bound and the dichotomy with the error \(p\,h(p)\) and has unbounded
degeneration, so a linear error is the threshold; and the tower of widths
\((k+1)!\) satisfies the Teichm\"uller contraction and the dichotomy with a
linear error but has uncapped jumps and unbounded degeneration. The widths of the
class of Problem~4.3 are the pseudo-Siegel degenerations: the quadratic-like
degeneration is unbounded in some classes of unbounded satellite type
(S{\o}rensen; Cheraghi--Shishikura), whose class of quadratic growth is a published
case. The conformal estimates behind the Gr\"otzsch step (4.6) and the Teichm\"uller
contraction are proved in \texttt{Mlc/Research/ExtremalLength.lean}: the series law,
the parallel law, \(W(G\setminus G')\ge\delta^{-1}W(G)\) and the pullback law
\(W(\Gamma_X)\le d\,W(\Gamma_Y)\). The abstract extremal length agrees with the classical
modulus on the model case: in the plane with Lebesgue measure, the radial segments of the
annulus \(\{r<|z|<R\}\) have extremal length \(\log(R/r)/2\pi\) and its separating circles
\(2\pi/\log(R/r)\) (\texttt{extLength\_radialFamily}, \texttt{extLength\_circleFamily} in
\texttt{Mlc/Research/AnnulusModulus.lean}). The same holds for a rectangle: the horizontal
segments of \((x_0,x_1]\times(y_0,y_1]\) have extremal length \((x_1-x_0)/(y_1-y_0)\)
(\texttt{extLength\_horizFamily} in \texttt{Mlc/Research/RectangleModulus.lean}), which
computes the widths of the rectangles of Figure~13. For the nested pieces of a
parameter neighbourhood the Gr\"otzsch--Teichm\"uller estimate is proved: if the closure of
the piece \(n\) of a nest in the disk \(B(c,R)\) has a point at distance \(\delta\) from \(c\),
the moduli of its first \(n\) annuli (the extremal widths of the \(C^1\) loops around the next
piece) add up to at most \(2\pi\log(4(R+\delta)/\delta)/(\log 2)^2\)
(\texttt{ConnectedNest.sum\_extWidth\_loops\_le} in \texttt{Mlc/Research/NestModulus.lean}),
so a nest whose moduli add up to infinity shrinks to \(c\), and Lemma~T follows from the
phase--parameter relation in moduli form (\texttt{Geometry.PhaseParameter}); the concentric
disks show that this hypothesis has instances. The three statements of Problem~4.3 and the a priori
bounds are invariant under the removal of a finite prefix of levels, which is the
untuning of a tuned parameter (\texttt{lemmas\_shift\_iff} in
\texttt{Mlc/Research/TailInvariance.lean}): they hold for a tower exactly when they hold
for the tower without its first \(n\) levels. This is the exact, combinatorial
self-similarity of \(M\); the uniform metric comparison of the tuned copies stays open.
The statements of the programme
are hypotheses; the roadmap lists their sources and an attack plan for each.

\section{Progress dashboard}

The baseline (update~0) is the state of the literature: MLC follows from Dudko's
Problems 4.3 and 4.4, which stay open. The area \(4\pi\) of the disk \(|c|\le 2\)
only normalizes the scale of the share below. The area of the parameters that the
published theorems leave open is unknown (expected to be small, without a
certificate), so the baseline row claims no unresolved area and the report states no
reduction factor relative to the literature. The first row of
Table~\ref{tab:commits} therefore shows the open problems instead of a
percentage.

The headline number is the certified resolved share
\[
  1-A_t=1-\frac{\mu^*(R_t)}{4\pi},
\]
where \(R_t\subseteq M\) is the unresolved set: the parameters without a
proved local-connectivity certificate at update \(t\). A parameter is resolved
when it lies outside \(M\) or in \(\operatorname{int}M\), where local
connectivity is automatic. The value \(100\%\) means \(\mu^*(R_t)=0\).
The share is a by-product of the kernel-checked grid: the cells that no
certificate removes form a neighbourhood of \(\partial M\) at the resolution of the
grid, which shrinks with the depth of the grid and does not measure progress on
Problems 4.3 and 4.4.
Table~\ref{tab:commits} lists every update after the baseline, together with
the resolved share and the result that raised it. Figure~\ref{fig:map} marks
on \(M\) the parameters that each update certified.

The updates are stored in the repository tree (\texttt{docs/progress}), numbered from
the baseline, and dated by the decision that made them; the table needs no Git history.

\begingroup\footnotesize
\begin{xltabular}{\linewidth}{@{}ll%
  >{\hsize=.8\hsize\raggedright\arraybackslash}X%
  >{\raggedleft\arraybackslash}p{23mm}%
  >{\hsize=1.2\hsize\raggedright\arraybackslash}X@{}}
  \caption{Updates after the literature baseline. The number counts the updates
  from the baseline (0), and the date is that of the decision. A coloured square
  marks an update with a layer on the map; a dash marks an update that resolves
  nothing beyond the literature. The resolved share is \(1-A_t\), where \(A_t\) is the
  certified upper bound on \(\mu^*(R_t)/(4\pi)\).}\label{tab:commits}\\
  \toprule
  \# & Date & Description & Resolved share & What it resolved \\
  \midrule
  \endfirsthead
  \toprule
  \# & Date & Description & Resolved share & What it resolved \\
  \midrule
  \endhead
  \bottomrule
  \endfoot
  \bottomrule
  \endlastfoot
  \input{docs/mlc_progress_commits.tex}
\end{xltabular}
\endgroup

\begin{figure}[ht]
  \centering
  \includegraphics[width=0.74\linewidth]{docs/mlc_progress_map.pdf}
  \caption{Parameter plane \(c=x+iy\). White outline: numerical rendering of
  the border of \(M\) (an illustration that the Lean proof does not use). A
  colour marks the parameters certified by the update with the same colour in
  Table~\ref{tab:commits}. The layers are stacked from the newest, largest
  region at the bottom to the oldest on top, so that every update stays
  visible; the dashed circle is the envelope \(|c|=2\).}
  \label{fig:map}
\end{figure}

\paragraph{What the value $100\%$ means.}
MLC means \(R=\varnothing\), hence \(A=0\) and the resolved share is
\(100\%\). The converse fails: a nonempty \(R\) of measure zero also gives
\(100\%\). Moreover \(R\subseteq\partial M\), because local connectivity is
automatic at interior points, and the area of \(\partial M\) is itself
unknown; MLC is not known to decide it, since a locally connected planar
continuum can have a boundary of positive area (an Osgood curve). Completion
is therefore tracked by the open root obligations together with the share.

The formalization-debt proxy is
\[
  D_{\mathrm{Lean}}
  =
  \frac{\text{number of open root obligations}}
       {1+\text{achieved milestone weight}}.
\]
It is an audit gate rather than a measure of the Mandelbrot residual set.
The commit hook rejects a proof-progress commit unless it improves the certified
resolved share, the share of classes with a Lean certificate, or this value, and
worsens none of them; it then records the update in the progress ledger.

\begin{table}[ht]
  \centering\footnotesize
  \input{docs/mlc_progress_routes.tex}
  \caption{Route-planning priors. Ranges and priors are planning estimates.}
  \label{tab:routes}
\end{table}

\section{Coverage-Uncertainty Profile}

Let \(R_t\subseteq\partial M\) be the parameters without a verified
local-connectivity certificate at stage \(t\). The area component is
\[
  A_t=\frac{\mu^*(R_t)}{4\pi}.
\]
The checked bound is
\[
  A_t\leq\frac{2\cdot 556432512}{4^{16}}\cdot\frac{1}{4\pi}\approx 0.0206.
\]
The actual value remains unknown. The first gate \(A_t\leq 0.05\) (a resolved
share of \(95\%\)) is reached. The next gate is \(A_t\leq 0.01\) (\(99\%\)):
\[
  \mu^*(R_t)\leq \frac{4\pi}{100}=\frac{\pi}{25},
  \qquad
  A_t\leq 0.01.
\]
The gates are milestones of the grid computation. The cells that no certificate removes
form a neighbourhood of \(\partial M\), whose area shrinks with the depth and says nothing
about the classes that the published theorems leave open; passing the next gate
requires a finer grid or a certified covering estimate of that neighbourhood.

Area alone cannot measure all progress. The report therefore uses the
\emph{Coverage-Uncertainty Profile}
\[
  (A_t,C_t,K_t,B_t),
\]
where:

\begin{itemize}
  \item \(C_t\) is a weighted count of parameter classes outside the verified
        coverage ledger;
  \item \(K_t\) is a weighted complexity term for unresolved return depth and
        relative period;
  \item \(B_t\) is a weighted deficit for missing quantitative separation
        budgets.
\end{itemize}

Each component is nonincreasing when a class certificate is retained. For a
catalogued class \(C_i\), completion strictly decreases the class component
by at least its assigned weight. The catalogue has eight classes of equal
weight, each with the statements of its chain; the proof of one statement lowers the
class component by the statement's share of every class that needs it. The profile
is a dashboard invariant; its components record how much of the parameter space
remains uncertain.

\section{Route priors}

The chart uses broad planning priors for four routes over a sustained
multi-year effort:

\begin{center}
\small
\setlength{\tabcolsep}{4pt}
\begin{tabular}{p{0.44\linewidth}cc}
\toprule
Route & Planned residual reduction & Success prior \\
\midrule
Problem 4.3: remaining unbounded satellite pseudo-Siegel bounds &
\(10\)--\(30\%\) & \(12\%\) \\
Problem 4.4: virtual bounded-period satellite interpolation &
\(24\)--\(46\%\) & \(20\%\) \\
Problem 4.4: virtual near-neutral interpolation &
\(18\)--\(42\%\) & \(12\%\) \\
Unified phase-parameter transfer conditional on genuine geometry &
\(38\)--\(72\%\) & \(45\%\) \\
\bottomrule
\end{tabular}
\end{center}

These ranges are planning estimates. Their purpose is
to prioritize theorem acquisition and estimate which route can settle the largest
part of the residual set, the parameters that the published theorems leave open.
The planned residual reduction is a relative research-management prior for that set
and is unrelated to the area of the grid. The ranges must be revised after every proved
class.

\section{Mathematical frontier}

\subsection{Combinatorial renormalization}

Formalize primitive and satellite copies, relative periods, the finite
chain \(M^{(0)}\supsetneq\cdots\supsetneq M^{(n)}\), \(1/2\)-decorations,
bounded-period profiles, virtual bounded-type profiles, and virtual
near-neutral profiles.

\subsection{Dynamical and geometric data}

Construct genuine small Julia sets, phase-indexed first-return domains,
critical-orbit windows, branch degree, distortion, collars, partial
invariance, and hulls inside \(J_c\). Prove the actual hull and return
estimates rather than storing them as abstract fields.

\subsection{Width and parameter transfer}

Formalize conformal modulus, Gr\"otzsch, quasi-additivity, the Wave Lemma,
and the Covering Lemma with all degree and topology-dependent thresholds.
Then construct moving parameter puzzles, holomorphic motion, winding-number
control, and a divergent parameter-side budget.

\subsection{Coverage theorem}

For the final candidate system \(A\), prove
\[
  \forall c\in M,\quad
  c\in\operatorname{int}M
  \;\lor\;
  A.\mathrm{Eligible}(c)
  \;\lor\;
  \operatorname{Nonempty}
  (\mathrm{VirtualMoleculeAPrioriControl}\ c).
\]
This is the first unconditional topological milestone. It still has to be
assembled with the concrete Problem 4.3 realization before full MLC follows.

\section{Completion criteria}

The next accepted progress commits must establish, in order:

\begin{enumerate}
  \item a strict increase of the certified resolved share (a strict decrease
        of the residual-area envelope) and a strict decrease of the
        formalization-debt proxy;
  \item an explicit class predicate independent of local connectedness;
  \item a concrete geometric witness for that class;
  \item a proved parameter realization and separation budget;
  \item a residual outer-measure update or another rigorous component of the
        Coverage-Uncertainty Profile;
  \item a strict graph, axiom, and build audit.
\end{enumerate}

The terminal commit must discharge both open root obligations and instantiate
`RootInput` without an additional input argument.

\end{document}
